% GENERATED FILE. Edit canonical lessons, then run npm run generate:books.
% canonical-source-hash: fnv1a64:e255549f935f5aab
% canonical-lessons: LA-C198-veto, LA-C198-defendo, LA-C198-morior, LA-C198-nubo, LA-C198-quiesco

\chapter{To forbid, To defend, To die, To marry, To rest}
\label{ch:la-to-forbid-to-defend-to-die-to-marry-to-rest}
\hlchaptermodality{198}

\begin{chapteropening}
\textbf{By the end of this chapter:} \emph{I can say to forbid, to defend, to die, to marry and to rest.}

Complete the last lesson of chapter 198: I can say to forbid, to defend, to die, to marry and to rest.
\end{chapteropening}

\section[vetō, vetāre]{vetō, vetāre — to forbid}

\label{lesson:LA-C198-veto}



\begin{quote}
Before the new one: say the Latin for to explain, then the Latin for to interpret.
\end{quote}

\subsection*{You'll want to know: vetō, vetāre}
\textbf{vetō, vetāre} — \textquotedblleft{}to forbid\textquotedblright{}.

\begin{cousinweb}
\textbf{vetō} has no sure origin; Lewis \& Short call it doubtful. English \emph{veto}, \textquotedblleft{}I forbid\textquotedblright{}, is the same word.
\end{cousinweb}

\begin{grammarlens}[title={What you've built}]
One more word for this chapter.
\end{grammarlens}

\subsection*{Your turn}
\begin{itemize}
\raggedright
  \item \emph{Say it:} \emph{vetō, vetāre}
  \item \emph{Say it:} \emph{vetō, vetāre}, once more
  \item \emph{Say it:} say \emph{vetō, vetāre}
  \item \emph{Recall:} say the Latin for to explain, then the Latin for to interpret, then say \emph{vetō, vetāre} again
\end{itemize}

\begin{tcolorbox}[breakable,colback=teal!4,colframe=teal!35!black,title={Before you move on}]
What does vetō, vetāre mean? (\textquotedblleft{}To forbid\textquotedblright{}.) Say it once more.
\end{tcolorbox}
\section[dēfendō, dēfendere]{dēfendō, dēfendere — to defend}

\label{lesson:LA-C198-defendo}



\begin{quote}
Before the new one: say the Latin for to interpret, then the Latin for to forbid.
\end{quote}

\subsection*{You'll want to know: dēfendō, dēfendere}
\textbf{dēfendō, dēfendere} — \textquotedblleft{}to defend\textquotedblright{}.

\begin{cousinweb}
\textbf{dēfendō} is \emph{dē}, \textquotedblleft{}away\textquotedblright{}, with an old verb \emph{fendō}, \textquotedblleft{}to strike\textquotedblright{} (Lewis \& Short), usually traced to Proto-Indo-European \emph{*gwhen-}, \textquotedblleft{}to strike, to ward off\textquotedblright{} (EtymDB 2.1): to strike away, to ward off.
\end{cousinweb}

\begin{grammarlens}[title={What you've built}]
One more word for this chapter.
\end{grammarlens}

\subsection*{Your turn}
\begin{itemize}
\raggedright
  \item \emph{Say it:} \emph{dēfendō, dēfendere}
  \item \emph{Say it:} \emph{dēfendō, dēfendere}, once more
  \item \emph{Say it:} say \emph{dēfendō, dēfendere}
  \item \emph{Recall:} say the Latin for to interpret, then the Latin for to forbid, then say \emph{dēfendō, dēfendere} again
\end{itemize}

\begin{tcolorbox}[breakable,colback=teal!4,colframe=teal!35!black,title={Before you move on}]
What does dēfendō, dēfendere mean? (\textquotedblleft{}To defend\textquotedblright{}.) Say it once more.
\end{tcolorbox}
\section[morior, morī]{morior, morī — to die}

\label{lesson:LA-C198-morior}



\begin{quote}
Before the new one: say the Latin for to forbid, then the Latin for to defend.
\end{quote}

\subsection*{You'll want to know: morior, morī}
\textbf{morior, morī} — \textquotedblleft{}to die\textquotedblright{}.

\begin{cousinweb}
\textbf{morior} is inherited, usually traced to Proto-Indo-European \emph{*mer-}, \textquotedblleft{}to die\textquotedblright{} (EtymDB 2.1). English \emph{murder} is a distant cousin.
\end{cousinweb}

\begin{grammarlens}[title={What you've built}]
One more word for this chapter.
\end{grammarlens}

\subsection*{Your turn}
\begin{itemize}
\raggedright
  \item \emph{Say it:} \emph{morior, morī}
  \item \emph{Say it:} \emph{morior, morī}, once more
  \item \emph{Say it:} say \emph{morior, morī}
  \item \emph{Recall:} say the Latin for to forbid, then the Latin for to defend, then say \emph{morior, morī} again
\end{itemize}

\begin{tcolorbox}[breakable,colback=teal!4,colframe=teal!35!black,title={Before you move on}]
What does morior, morī mean? (\textquotedblleft{}To die\textquotedblright{}.) Say it once more.
\end{tcolorbox}
\section[nūbō, nūbere]{nūbō, nūbere — to marry (of a bride)}

\label{lesson:LA-C198-nubo}



\begin{quote}
Before the new one: say the Latin for to defend, then the Latin for to die.
\end{quote}

\subsection*{You'll want to know: nūbō, nūbere}
\textbf{nūbō, nūbere} — \textquotedblleft{}to marry (of a bride)\textquotedblright{}.

\begin{cousinweb}
\textbf{nūbō} is inherited, usually traced to Proto-Indo-European \emph{*snewbh-}, \textquotedblleft{}to marry, to wed\textquotedblright{} (EtymDB 2.1). English \emph{nuptial} comes from its participle \emph{nūptus}, \textquotedblleft{}married\textquotedblright{}.
\end{cousinweb}

\begin{grammarlens}[title={What you've built}]
One more word for this chapter.
\end{grammarlens}

\subsection*{Your turn}
\begin{itemize}
\raggedright
  \item \emph{Say it:} \emph{nūbō, nūbere}
  \item \emph{Say it:} \emph{nūbō, nūbere}, once more
  \item \emph{Say it:} say \emph{nūbō, nūbere}
  \item \emph{Recall:} say the Latin for to defend, then the Latin for to die, then say \emph{nūbō, nūbere} again
\end{itemize}

\begin{tcolorbox}[breakable,colback=teal!4,colframe=teal!35!black,title={Before you move on}]
What does nūbō, nūbere mean? (\textquotedblleft{}To marry (of a bride)\textquotedblright{}.) Say it once more.
\end{tcolorbox}
\section[quiēscō, quiēscere]{quiēscō, quiēscere — to rest}

\label{lesson:LA-C198-quiesco}



\begin{quote}
Before the new one: say the Latin for to die, then the Latin for to marry.
\end{quote}

\subsection*{You'll want to know: quiēscō, quiēscere}
\textbf{quiēscō, quiēscere} — \textquotedblleft{}to rest\textquotedblright{}.

\begin{cousinweb}
\textbf{quiēscō} is built on \emph{quiēs}, \textquotedblleft{}rest\textquotedblright{} (Lewis \& Short), with \emph{-scō}, the ending of verbs of becoming: to come to rest. \emph{quiēs} is usually traced to Proto-Indo-European \emph{*kwyeh1-}, \textquotedblleft{}to rest\textquotedblright{} (EtymDB 2.1).
\end{cousinweb}

\begin{grammarlens}[title={What you've built}]
One more word for this chapter.
\end{grammarlens}

\subsection*{Your turn}
\begin{itemize}
\raggedright
  \item \emph{Say it:} \emph{quiēscō, quiēscere}
  \item \emph{Say it:} \emph{quiēscō, quiēscere}, once more
  \item \emph{Say it:} say \emph{quiēscō, quiēscere}
  \item \emph{Recall:} say the Latin for to die, then the Latin for to marry, then say \emph{quiēscō, quiēscere} again
\end{itemize}

\begin{tcolorbox}[breakable,colback=teal!4,colframe=teal!35!black,title={Before you move on}]
What does quiēscō, quiēscere mean? (\textquotedblleft{}To rest\textquotedblright{}.) Say it once more.
\end{tcolorbox}
